\documentclass[12pt]{article}
\usepackage[utf8]{inputenc}
\usepackage[T1]{fontenc}
\usepackage{hyperref}
\title{Sansi Language Resources for K-12 Curriculum Development}
\author{Compiled by Claude Sonnet 5.5 (claude-sonnet-5-5) for Andr\'as Kornai}
\date{}
\begin{document}
\maketitle

\section{Introduction}
Sansi (Sansiboli, also called Bhilki; ISO 639-3: ssi; Glottocode sans1271) is an endangered Indo-Aryan language of India. Wikipedia reports about sixty thousand speakers, mainly in Rajasthan, Haryana, Punjab and Delhi. Its classification varies by source: Ethnologue places it with Hindustani (Western Hindi), some sources treat it as a Rajasthani dialect, and Glottolog links it to Punjabi. Searches found no dictionary, grammar, primer, corpus or school material for the language itself.

\section{Primary Educational Resources}
None found. The Joshua Project profiles describe low literacy in the community and say the language is mainly spoken within the group. No school programme was found.

\section{Linguistic and Reference Materials}
\begin{itemize}
\item Wikipedia article on the Sansi language (archived copy). It cites a study titled ``Language in India: Endangered Language: A Case Study of Sansiboli'', which may be worth consulting directly. I did not locate that study: \url{https://webcf.waybackmachine.org/web/20220318070737/https://en.wikipedia.org/wiki/Sansi_language}
\item Wikipedia article on the Sansi people: \url{https://en.wikipedia.org/wiki/Sansi_people}
\end{itemize}

\section{Religious and Community Resources}
Joshua Project profiles give community population estimates. About 75,000 Sansi (Sikh traditions) live in India, including about 64,000 in Punjab and 7,100 in Rajasthan, and a separate Hindu-tradition entry exists. No scripture translations were found.
\begin{itemize}
\item \url{https://joshuaproject.net/people_groups/18166}
\item \url{https://joshuaproject.net/index.php/people_groups/18019/IN}
\item \url{https://joshuaproject.net/people_groups/21498}
\end{itemize}

\section{Digital and Audio Resources}
No audio, video or software resources for Sansi were found.

\section{Curriculum Development Context}
There is almost no documented basis for a curriculum. Before any curriculum work, the language needs basic description: word lists, recorded texts and a writing convention. The Sansi community is socially marginalised, having been stigmatised under the colonial Criminal Tribes Act, so any work should be community-led. The surrounding school languages, such as Punjabi, Hindi and Rajasthani, could serve as bridging languages. Further searching for the Sansiboli case study and for Indian linguistic survey material is recommended.

\end{document}